%% kapa: exam=1 round=24 year=2013 period=1 form=A subject=civil_law subject_ko=민법 seq=32 number=32
%% source: Q-Net 감정평가사 기출문제 (cst00309) artlSeq=5073802 file=extracted/1008/5073802/2013년 제24회 감정평가사 1차시험 기출문제/1교시(A형).hwp
%% conversion: kordoc 4.24.51 -> tools/gen1.py (auto); review=proofread-sonnet; flags=none
\begin{question}{32}
관습법상의 법정지상권에 관한 설명으로 옳지 않은 것은? (다툼이 있으면 판례에 의함)
\begin{choices}
  \item 토지의 매수인이 소유권이전등기를 하기 전에 토지소유자의 승낙을 받아 건물을 신축했으나 매매계약이 해제된 경우 관습법상의 법정지상권이 성립하지 않는다.
  \item 관습법상 법정지상권을 미리 포기하기로 하는 약정은 효력이 있다.
  \item 토지공유자의 한 사람이 과반수 지분을 갖는 일부 공유자만의 동의를 얻어 건물을 신축한 후 이를 타인에게 양도한 경우 관습법상의 법정지상권이 성립하지 않는다.
  \item 토지 소유자가 담보의 목적으로 채권자에게 그 토지위에 가등기를 마친 뒤 건물을 신축하였고 그 후 채권자가 가등기에 기한 본등기를 한 경우 특별한 사정이 없으면 관습법상 법정지상권이 성립한다.
  \item 토지를 제3자에게 매도하고 토지소유권이전등기를 하기 전에 매도인이 건물을 신축하고 그 후에 토지소유권이전등기를 한 경우 특별한 사정이 없으면 관습법상 법정지상권이 성립하지 않는다.
\end{choices}

eviewflag{choices 3-4 overlap in rhwp render; only part verified}
\end{question}
\answer{4}
